\ssszh $\DGCat$ 的一个很有用的性质是，可以把余完备的 dg 范畴张量起来
%
\[
   \sC \text{ 与 } \sD \rightsquigarrow \sC \otimes \sD.
\]
%
更确切地说，$\DGCat$ 关于 Lurie 张量积是对称幺半 $\infty$-范畴。

在实际的许多例子中，下面这种可以动手算的描述很有用。

\begin{ex}设 $\sC$ 与 $\sD$ 是紧生成的，比如分别由对象集 $c_\alpha$ 与 $d_\beta$ 紧生成。则张量积由简单张量 $c_\alpha \otimes d_\beta$ 构成的集合紧生成，其同态由 K\"unneth 公式给出，即
%
\[
\Hom( c_\alpha \otimes d_\beta, c_{\alpha'} \otimes d_{\beta'}) \simeq
\Hom(c_\alpha, c_{\alpha'}) \otimes \Hom(d_\beta, d_{\beta'}).
\]
%
\end{ex}

上述例子有一些特殊情形，希望它们能让读者感受到，这个张量积与读者从函数-层对应中得来的直觉是一致的，参见第 \ref{ss:opsfinitesets} 小节。

\begin{ex} 给定代数簇 $X$ 与 $Y$，拟凝聚层范畴之间有典范等价
%
\[
   \on{QCoh}(X) \otimes \on{QCoh}(Y) \simeq \on{QCoh}(X \times Y).
\]
%
具体地说，给定 $\on{QCoh}(X)$ 的对象 $\sM$ 与 $\on{QCoh}(Y)$ 的对象 $\sN$，左边的简单张量 $\sM \otimes \sN$ 被送到外积对象 $$\sM \boxtimes \sN
\simeq \pi_X^* \sM
\overset{}\otimes \pi_Y^* \sN.$$
\end{ex}

\begin{ex} 仍取上例中的 $X$ 与 $Y$。D-模范畴之间也有典范等价
%
\[
     \Dmod(X) \otimes \Dmod(Y) \simeq \Dmod(X \times Y).
\]
%
同样地，给定 $\Dmod(X)$ 的对象 $\sM$ 与 $\Dmod(Y)$ 的对象 $\sN$，左边的简单张量 $\sM \otimes \sN$ 被送到外积对象
%
\[
      \sM \boxtimes \sN \simeq \pi_X^! \sN \overset{!}{\otimes} \pi_Y^! \sM.
\]
%
\zy{[译者按（上下文推断，无外部依据）：原文此处把 $\sM$ 与 $\sN$ 的位置写颠倒了。按上例的样式，应为 $\pi_X^! \sM
\overset{!}{\otimes} \pi_Y^! \sN$；原文写作 $\pi_X^! \sN \overset{!}{\otimes} \pi_Y^! \sM$。本译本照原文排印，读者按此理解即可。]}
\end{ex}

\begin{ex}给定 $k$-代数 $A$ 与 $B$，有典范等价
%
\[
    A\mod \otimes \hspace{.2mm} B\mod \simeq (A \underset{k} \otimes
    B)\mod.
\]
%
注意这推广了前两个例子的仿射情形。
\end{ex}

需要一提的是，这个幺半结构是幺半的，其单位是向量空间范畴 $\Vect$。特别地，对任意 dg 范畴 $\sC$，有典范等价
%
\[
   \Vect \otimes \hspace{.7mm}  \sC \simeq \sC.
\]
%
读者不妨在 $\sC$ 紧生成的情形直接验证这一点。

\ssszh 给定 dg 范畴 $\sC,\sD$，第 \ref{s:homdgcats} 小节中的范畴 $\udHom(\sC,
\sD)$ 关于 $\DGCat$ 上的幺半结构是内 Hom 对象。

这里有一个技术要点：它并不等同于把 $\sC$ 与 $\sD$ 看成 $\DGCat$ 的对象之后二者之间的态射范畴。准确地说，态射范畴是把 $\udHom(\sC, \sD)$ 截断到 $\infty$-群胚，也就是丢掉不可逆的自然变换。人们预期 $\DGCat$ 有典范的提升，成为对称幺半 $(\infty,2)$-范畴，其中同态确实就是 $\udHom(-, -)$；但据我们所知，文献中目前还没有这一构造。

\ssszh 借助幺半结构，可以谈论 $\DGCat$ 的代数对象，即幺半 dg 范畴。粗略地说，它由一个 dg 范畴 $\sA$ 连同乘法
%
\[
    \sA \otimes \sA \rightarrow \sA,
\]
%
组成，该乘法在相容同伦的意义下结合。我们强调，相容同伦是额外的结构，是幺半 dg 范畴完整数据的一部分。

下面是幺半 dg 范畴的一些基本例子。

\begin{ex}\label{ex:qcohalg} 对代数簇 $X$，$\QCoh(X)$ 在 $*$-张量积下是幺半 dg 范畴，$\Dmod(X)$ 在 $!$-张量积下是幺半 dg 范畴。

给定映射 $f: X_1 \rightarrow X_2$，相应的拉回函子
%
\[
f^*: \QCoh(X_2) \rightarrow \QCoh(X_1) \quad \text{与} \quad f^!: \Dmod(X_2)
\rightarrow \Dmod(X_1)
\]
%
带有典范的幺半结构数据。
\end{ex}

\begin{ex} 设簇 $Y$ 是簇范畴中的代数对象。也就是说，$Y$ 带有映射
%
\[
    \mu: Y \times Y \rightarrow Y,
\]
%
它在显然的意义下结合。此时 $\QCoh(Y)$ 与 $\Dmod(Y)$ 都获得「卷积」幺半结构，由沿 $\mu$ 的 $*$-推前给出。

给定代数的映射 $f:Y_1 \rightarrow Y_2$，相应的推前映射
%
\[
 f_*: \QCoh(Y_1) \rightarrow \QCoh(Y_2) \quad \text{与} \quad f_*: \Dmod(Y_1)
 \rightarrow \Dmod(Y_2)
\]
%
带有典范的幺半结构数据。

我们主要关心的是 $Y = K$ 为代数群、$\mu$ 为其乘法这一情形，由此得到先前见过的群代数 $\Dmod(K)$（更确切地说，是它的典范 dg 提升）。

类似地，给定群的同态 $\phi: K_1 \rightarrow K_2$，得到群代数的同态
%
\[
   \phi_*: \Dmod(K_1) \rightarrow \Dmod(K_2).
\]
%
\end{ex}

\begin{re} 把前两段中的 $\Dmod(K)$ 换成 $\QCoh(K)$，可以得到另一个卷积下的幺半 dg 范畴。$\Dmod(K)$ 的模有时称作强范畴表示，$\QCoh(K)$ 的模则称作弱范畴表示。

有一个或许能澄清问题的注：给定从簇到幺半 $\infty$-范畴 $(\sS, \otimes)$ 的任意幺半函子，都可得到类似的、把 $K$ 线性化的卷积代数。特别地，可以谈论 $\on{Fun}(K(\Fq))$ 的更高范畴版本，例如把层换成 $K$ 上层的范畴。\end{re}

\ssszh 给定幺半 dg 范畴 $\sA$，有与之相伴的左模与右模 $\infty$-范畴，分别记作
%
\[
  \sA\mod \quad \text{与} \quad \rmod\sA.
\]
%
说得不那么正式，一个左模由一个余完备 dg 范畴 $\sM$ 连同作用映射
%
\[
  \sA \otimes \sM \rightarrow \sM
\]
%
组成，该作用在相容同伦的意义下「结合」；右模类似。同样，相容同伦是指定模的数据的一部分。

给定两个左模 $\sM_1$ 与 $\sM_2$，有 $\sA$-等变函子构成的 dg 范畴
%
\[
   \Hom_{\sA\mod}(\sM_1, \sM_2)
\]
%
（$\sA$ 上右模的情形类似）。具体地说，它可以借助 bar 复形、作为 dg 范畴半余单纯图的极限来计算
%
\begin{equation} \label{e:hombar}
   \scriptsize
   \xymatrix{  \udHom(\sM_1, \sM_2)
   \ar@<-.4ex>[r] \ar@<.4ex>[r] &  \hspace{.5mm} \udHom(\sA \otimes \sM_1,
   \sM_2)
   \ar@<-.8ex>[r] \ar@<.8ex>[r] \ar[r] & \hspace{.5mm} \udHom(\sA^{\otimes 2}
   \otimes \sM_1, \sM_2)   \hspace{.5mm} \cdots \hspace{.5mm},
   }
   \end{equation}
%
其中各箭头按 bar 消解的标准模式，由 $\sA$ 在 $\sM_1$ 与 $\sM_2$ 上的作用映射、以及乘法映射 $\sA^{\otimes 2} \rightarrow \sA$ 诱导。

此外，给定右模 $\sM$ 与左模 $\sN$，可以作它们的张量积
%
\[
     \sM  \underset{\sA}{\otimes} \sN.
\]
%
与同态的情形一样，它可以借助 bar 消解、作为余完备 dg 范畴半单纯图的余极限来计算
\zy{[译者按（上下文推断，无外部依据）：下式第二项中的 $N$，按全式其余各项的写法应作花体的 $\sN$。原文此处漏写花体，本译本照原文排印。]}
%
\begin{equation} \label{e:tensorbar}
\xymatrix{  \cdots \hspace{.5mm} \sM \otimes \sA^{\otimes 2} \otimes N
\ar@<-.8ex>[r] \ar@<.8ex>[r] \ar[r]   &  \hspace{.5mm} \sM \otimes \sA \otimes
\sN \ar@<-.4ex>[r]
\ar@<.4ex>[r]  &  \hspace{.5mm} \sM \otimes  \sN}.
\end{equation}
%


现在给出几个基本例子。

\begin{ex} 对任意左 $\sA$-模 $\sM$，有显然的等价
%
\[
   \on{Hom}_{\sA\mod}(\sA, \sM) \simeq \sM \simeq \sA \underset \sA \otimes
   \sM
\]
%
右模的情形类似。
\end{ex}

\begin{ex} \label{e:fibprodvar}给定簇的映射 $X \rightarrow Z
\leftarrow Y$，如例 \ref{ex:qcohalg} 那样，得到 $\QCoh(Z)$ 在 $\QCoh(X)$ 与 $\QCoh(Y)$ 上的典范作用。于是它们的张量积由下式给出
%
\[
    \QCoh(X)
    \underset{\QCoh(Z)}{\otimes} \QCoh(Y) \simeq \QCoh( X \underset{Z}{\times}
    Y),
\]
%
其中上述纤维积是导出的，即得到导出概形。\footnote{直白地说，这是逐仿射地用导出环 $\OO_X
\underset{\OO_Z}{\overset{L}{\otimes}} \OO_Y$ 粘起来的，与通常纤维积的情形完全一样。人们把结构层的负上同调群想成给簇添上「幂零绒毛」，就像一般概形的非既约性那样，只不过后者或许更为读者所熟悉。本书会遇到的一切导出概形都以这种方式产生。}
具体地说，这个等价把简单张量 $\sM \otimes \sN$ 送到它们的外积在 $X \underset{Z}{\times} Y$ 上的限制。

类似地，它们的同态也由下式给出
%
\[
 \Hom_{\QCoh(Z)}(\QCoh(X), \QCoh(Y)) \simeq \QCoh( X \underset{Z}{\times}
 Y).
\]
%
具体地说，把纤维积上的拟凝聚层 $\sS$ 对应到积分变换
%
\[
  \QCoh(X) \xrightarrow{ \pi_X^*} \QCoh(X \underset Z \times Y) \xrightarrow{ -
  \overset * \otimes \sS} \QCoh(X \underset Z \times Y) \xrightarrow{\pi_{Z,
  *}} \QCoh(Y).
\]
%

对 D-模，同样的断言作必要的调整后成立。也就是说，有等价
%
\begin{align*}
  \Dmod(X) \underset{\Dmod(Z)}{\otimes} \Dmod(Y) &\simeq \Dmod(X \overset{\on{cl}}{{\underset Z {\times}}} Y) \text{，以及}  \\
  \Hom_{\Dmod(Z)}(\Dmod(X), \Dmod(Y)) &\simeq \Dmod(X \overset{\on{cl}}{{\underset Z {\times}}} Y),
\end{align*}
%
它们同样由外张量积与积分变换给出。这里上标 $\on{cl}$ 表示经典的，即非导出的纤维积。但要注意，当导出纤维积的结构层只有有限多个非零上同调时（例如 $Z$ 光滑时），由于 D-模范畴在幂零加厚下不变，取簇的导出纤维积还是经典纤维积并不重要。

与我们已经遇到的其它层论陈述一样，这些等式是下面这个函数等式在层论中的类似物。给定有限集 $S$，把 $\Fun(S)$ 看作逐点乘法下的代数。给定有限集 $X \rightarrow Z
\leftarrow Y$，便有向量空间的典范同构
%
\[
    \Fun(X) \underset{\Fun(Z)}{\otimes} \Fun(Y) \simeq
    \Hom_{\Fun(Z)}(\Fun(X), \Fun(Y)) \simeq \Fun(X \underset Z \times Y).
\]
%
\end{ex}
